% Spine note preamble — not the Kingdom Come book preamble.
\usepackage[T1]{fontenc}
\usepackage[utf8]{inputenc}
\usepackage{lmodern}
\usepackage{microtype}
\usepackage{setspace}
\setstretch{1.28}

\usepackage[
  a4paper,
  margin=2.4cm,
  headheight=15pt,
  heightrounded
]{geometry}

\usepackage{amsmath,amssymb,amsthm,mathtools}
\usepackage{bm}

\usepackage{graphicx}
\usepackage{booktabs}
\usepackage{array}
\usepackage{longtable}
\usepackage{multirow}
\usepackage{tabularx}
\newcolumntype{Y}{>{\raggedright\arraybackslash}X}

\usepackage{float}
\usepackage{caption}
\captionsetup{font=small,labelfont=bf,format=hang,skip=6pt}

\usepackage{enumitem}
\setlist{nosep,leftmargin=*}

\usepackage{listings}
\usepackage{xcolor}
\definecolor{codebg}{RGB}{248,249,250}
\definecolor{codeframe}{RGB}{200,200,200}
\definecolor{codekw}{RGB}{0,70,140}
\lstdefinestyle{qga}{
  basicstyle=\ttfamily\footnotesize,
  backgroundcolor=\color{codebg},
  frame=single,
  rulecolor=\color{codeframe},
  breaklines=true,
  breakatwhitespace=true,
  columns=fullflexible,
  keepspaces=true,
  showstringspaces=false,
  tabsize=2,
  aboveskip=0.8em,
  belowskip=0.8em,
  keywordstyle=\color{codekw}\bfseries,
  commentstyle=\color{gray}\itshape,
  stringstyle=\color{teal},
}
\lstset{style=qga}

\usepackage[
  colorlinks=true,
  linkcolor=black!70!blue,
  citecolor=black!70!blue,
  urlcolor=black!60!blue,
  bookmarks=true,
  bookmarksnumbered=true,
  pdfstartview=FitH,
  pdftitle={QGA Spine Note: Quaternion Orders, Hopf, and the Hatcher Lift},
  pdfauthor={Aaron Michael Kinder}
]{hyperref}
\usepackage{xurl}
\urlstyle{tt}
\usepackage[nameinlink,noabbrev]{cleveref}

\usepackage{fancyhdr}
\pagestyle{fancy}
\fancyhf{}
\fancyhead[L]{\nouppercase{\leftmark}}
\fancyhead[R]{\thepage}
\renewcommand{\headrulewidth}{0.3pt}

\theoremstyle{definition}
\newtheorem*{qgatheorem}{Theorem}
\newtheorem*{qgamodel}{Model}
\newtheorem*{qgahypothesis}{Hypothesis}
\newtheorem*{qgasoftware}{Software fact}

\newcommand{\HH}{\mathbb{H}}
\newcommand{\RR}{\mathbb{R}}
\newcommand{\QQ}{\mathbb{Q}}
\newcommand{\ZZ}{\mathbb{Z}}
\newcommand{\CC}{\mathbb{C}}
